% =============================================================
%  Company Playbook — Đề án thành lập công ty
%  Tài liệu nền để 3 đồng sáng lập đối thoại và đồng thuận
%  Đà Nẵng, 2026 — Internal draft
%  Biên dịch: xelatex company-playbook.tex
% =============================================================
\documentclass[11pt,a4paper]{article}

\usepackage[margin=2.2cm,top=2cm,bottom=2.2cm]{geometry}
\usepackage{fontspec}
\usepackage{tikz}
\usepackage{xcolor}
\usepackage{tabularx}
\usepackage{booktabs}
\usepackage{enumitem}
\usepackage{titlesec}
\usepackage{fancyhdr}
\usepackage{hyperref}
\usepackage{setspace}
\usepackage{tcolorbox}
\tcbuselibrary{skins, breakable}
\usetikzlibrary{positioning, shapes.geometric, shapes.misc, arrows.meta,
                fit, calc, decorations.pathreplacing, backgrounds}

\setmainfont{Latin Modern Roman}
\setsansfont{Latin Modern Sans}

% ---------- colors ----------
\definecolor{ink}{HTML}{14213d}
\definecolor{accent}{HTML}{fca311}
\definecolor{mid}{HTML}{8d99ae}
\definecolor{hao}{HTML}{b08968}       % Thổ
\definecolor{phu}{HTML}{6c757d}       % Kim
\definecolor{hoang}{HTML}{2a9d8f}     % Mộc
\definecolor{warm}{HTML}{c44536}
\definecolor{soft}{HTML}{f4f1de}
\definecolor{leaf}{HTML}{3d5a80}

\hypersetup{colorlinks=true, linkcolor=ink, urlcolor=leaf}

% ---------- section styles ----------
\titleformat{\section}{\sffamily\bfseries\LARGE\color{ink}}{Chương \thesection.}{0.5em}{}
\titleformat{\subsection}{\sffamily\bfseries\large\color{ink}}{\thesubsection}{0.5em}{}
\titleformat{\subsubsection}{\sffamily\bfseries\normalsize\color{ink}}{}{0em}{}
\titlespacing*{\section}{0pt}{1.4em}{0.6em}
\titlespacing*{\subsection}{0pt}{1.0em}{0.3em}
\titlespacing*{\subsubsection}{0pt}{0.6em}{0.2em}

\setlist{nosep, leftmargin=1.3em, topsep=3pt, itemsep=2pt}
\setlength{\parskip}{4pt}
\setlength{\parindent}{0pt}
\onehalfspacing

% ---------- header/footer ----------
\pagestyle{fancy}
\fancyhf{}
\fancyhead[L]{\sffamily\footnotesize\color{mid} Company Playbook}
\fancyhead[R]{\sffamily\footnotesize\color{mid} Internal draft \textbullet{} 2026}
\fancyfoot[C]{\sffamily\footnotesize\color{mid}\thepage}
\renewcommand{\headrulewidth}{0pt}

% ---------- boxes ----------
\newtcolorbox{keybox}[1][]{
  enhanced, colback=soft, colframe=accent, boxrule=0.5pt,
  left=8pt, right=8pt, top=5pt, bottom=5pt, arc=1pt,
  fonttitle=\sffamily\bfseries\color{ink},
  #1
}
\newtcolorbox{warnbox}[1][]{
  enhanced, colback=warm!8, colframe=warm, boxrule=0.5pt,
  left=8pt, right=8pt, top=5pt, bottom=5pt, arc=1pt,
  fonttitle=\sffamily\bfseries\color{warm},
  #1
}
\newtcolorbox{quotebox}[1][]{
  enhanced, colback=white, colframe=ink!30, boxrule=0pt,
  borderline west={2pt}{0pt}{ink},
  left=10pt, right=8pt, top=4pt, bottom=4pt, arc=0pt,
  #1
}

\begin{document}

% =============================================================
% TITLE PAGE
% =============================================================
\thispagestyle{empty}
\vspace*{3cm}
\begin{center}
{\sffamily\bfseries\fontsize{28}{34}\selectfont\color{ink} Company Playbook}\\[0.5em]
{\sffamily\Large\color{mid} Đề án thành lập công ty — Bản nền để 3 người đồng thuận}\\[3em]

\begin{tikzpicture}[font=\sffamily]
  \node[circle, draw=ink, thick, fill=hao!30, minimum size=22mm] (h) at (0,0) {\textbf{Hảo}};
  \node[circle, draw=ink, thick, fill=phu!40, minimum size=22mm] (p) at (4,0) {\textbf{Phú}};
  \node[circle, draw=ink, thick, fill=hoang!30, minimum size=22mm] (ho) at (8,0) {\textbf{Hoàng}};
  \node[below=2mm of h, font=\footnotesize] {Strategy};
  \node[below=2mm of p, font=\footnotesize] {Platform};
  \node[below=2mm of ho, font=\footnotesize] {GTM};
  \draw[-{Latex[length=2.5mm]}, thick, hao] (h) -- node[above, font=\scriptsize, hao]{nuôi} (p);
  \draw[-{Latex[length=2.5mm]}, thick, hoang] (ho) to[bend left=40] node[above, font=\scriptsize, hoang]{lan toả} (h);
  \draw[-{Latex[length=2.5mm]}, thick, warm, densely dashed] (p) to[bend right=30] node[below, font=\scriptsize, warm]{cân} (ho);
\end{tikzpicture}

\vspace{3em}
{\large\color{ink}\sffamily
Đỗ Phúc Hảo \textbullet{} Lê Hữu Phú \textbullet{} Nguyễn Ngọc Hoàng
}\\[1em]
{\sffamily\color{mid} Đà Nẵng, tháng 4 năm 2026}
\end{center}
\vfill
\begin{center}
{\footnotesize\itshape\color{mid}
Tài liệu này là bản nháp nội bộ. Mục đích: làm nền để 3 người đối thoại, thống nhất tầm\\
nhìn, vai trò, lộ trình và các ranh giới — trước khi chính thức pháp lý hoá công ty.}
\end{center}
\clearpage

% =============================================================
% TABLE OF CONTENTS
% =============================================================
\tableofcontents
\clearpage

% =============================================================
% EXECUTIVE SUMMARY
% =============================================================
\section*{Tóm tắt điều hành}
\addcontentsline{toc}{section}{Tóm tắt điều hành}

\begin{keybox}[title=Một câu duy nhất mô tả công ty]
Chúng tôi xây \textbf{hệ thống AI Agent tự vận hành, có kiểm soát}, cho các
workflow tri thức, bắt đầu từ giáo dục đại học Việt Nam, với mục tiêu trở thành
công ty dẫn đầu Đông Nam Á về \textit{Agent Systems cho các ngành dịch vụ tri thức}
trong vòng 5 năm.
\end{keybox}

\textbf{Đích đến.} Trong 5 năm: một công ty 30--50 người, doanh thu dương, phục vụ
50--100 khách hàng giáo dục và dịch vụ tại Việt Nam \& Đông Nam Á, với một lõi
framework được tái sử dụng qua mọi dự án.

\textbf{Hành trình 3 năm:}
\begin{itemize}
\item \textbf{Năm 1 — Chứng minh.} 1 vertical (giáo dục), 2 use case (Admissions +
Internal Knowledge), 2--3 design partner trả phí, MVP chạy thật trong trường thật.
\item \textbf{Năm 2 — Nhân rộng.} 10--15 khách hàng giáo dục, solution đã được
chuẩn hóa, framework ổn định, bắt đầu thử vertical 2 (du lịch / dịch vụ).
\item \textbf{Năm 3 — Platform hóa.} Subscription model khi workflow đã lặp lại,
mở rộng địa bàn, bắt đầu tuyển thêm người.
\end{itemize}

\textbf{Con người.} Ba đồng sáng lập bổ sung nhau về kỹ năng lẫn tính cách:
Hảo (chiến lược, network học thuật), Phú (kỹ thuật, học trò cũ của Hảo), Hoàng
(kinh nghiệm thương mại, mạng lưới doanh nghiệp). Phân vai rõ ràng, có cơ chế
ra quyết định chung, có vesting 4 năm.

\textbf{Nguyên tắc cốt lõi.}
\begin{itemize}
\item \textbf{Solution ngoài, framework trong.} Thị trường trả tiền cho kết quả, không trả cho kiến trúc.
\item \textbf{Đơn giản trước, phức tạp sau.} Single agent tốt $>$ multi-agent không cần thiết.
\item \textbf{Khách trước, code sau.} Không xây thứ chưa có người trả tiền.
\item \textbf{Ba chữ ký cho mọi cam kết.} Không ai đơn phương hứa hẹn hay xây dựng.
\end{itemize}
\clearpage

% =============================================================
% CHAPTER I — BỐI CẢNH & CƠ HỘI
% =============================================================
\section{Bối cảnh \& Cơ hội}

\subsection{Thế giới đang ở đâu với Agentic AI (2025--2026)}

Ba năm từ 2023 đến 2026 là giai đoạn chuyển pha của AI doanh nghiệp. Giai đoạn
\textit{chatbot} (2023) đã qua. Giai đoạn \textit{copilot} (2024) đang bão hòa.
Giai đoạn hiện tại là \textbf{agentic} — tức là AI không chỉ gợi ý mà \textbf{tự vận
hành một phần workflow}, có công cụ, có quyền hành động, có đánh giá, có escalation.

Các tín hiệu quan trọng:
\begin{itemize}
\item \textbf{OpenAI \& Anthropic} đều công khai khuyến nghị: bắt đầu từ workflow
rõ ràng và guardrails rõ ràng, chỉ leo lên multi-agent khi thật sự cần.
\item \textbf{McKinsey (2025)} ước tính giá trị kinh tế từ agentic AI trên 4.400 tỷ
USD/năm toàn cầu, nhưng nhấn mạnh: giá trị đến từ tái thiết kế workflow, không phải
từ việc ``dán AI lên quy trình cũ''.
\item \textbf{Deloitte (2025)} xác định 5 vùng tác động cao nhất: customer support,
knowledge management, R\&D, supply chain, cybersecurity — \textbf{tất cả đều là
workflow tri thức}.
\item \textbf{Chi phí inference giảm 10--50$\times$} mỗi năm (Claude, GPT, Gemini,
và các model open-source). Rào cản kinh tế gần như không còn.
\end{itemize}

\begin{keybox}[title=Hàm ý cho chúng ta]
Cửa sổ thời gian cho công ty \textit{mới} xây solution vertical-specific
trong ngành tri thức đang mở. Nhưng cửa sổ này có thể đóng trong 2--3 năm khi
các big tech tung ra giải pháp vertical, và các công ty ``first movers'' đã
chiếm được design partner then chốt.
\end{keybox}

\subsection{Việt Nam: một thị trường ``dưới ngưỡng''}

Việt Nam có đặc điểm: doanh nghiệp và tổ chức giáo dục \textbf{đang dùng rất ít
AI thực sự}, dù họ đọc và nghe rất nhiều về AI. Lý do chính không phải kỹ thuật,
mà là:

\begin{itemize}
\item \textbf{Thiếu người Việt hiểu đồng thời} vừa sâu về công nghệ AI, vừa đủ
kinh nghiệm quy trình nghiệp vụ trong ngành.
\item \textbf{Các công ty công nghệ lớn} (FPT, VNG, CMC) tập trung vào tập đoàn
lớn, bỏ qua phân khúc doanh nghiệp vừa \& nhỏ và các tổ chức giáo dục tầm trung.
\item \textbf{Các startup AI Việt} đa số đuổi theo chatbot tiếng Việt, ít ai
hiểu sâu về agent systems.
\item \textbf{Nhà trường, công ty du lịch, tổ chức dịch vụ tri thức} có nhiều
quy trình lặp đi lặp lại, tài liệu rải rác, nhưng không có ai \textbf{đến gõ
cửa và hiểu họ}.
\end{itemize}

Đây là một thị trường ``dưới ngưỡng'' — nhu cầu có, chi trả được, nhưng
\textbf{chưa có ai đứng đúng chỗ để phục vụ}.

\subsection{Đà Nẵng: vị trí địa lý có chủ ý}

Đặt công ty tại Đà Nẵng không phải ngẫu nhiên. Đà Nẵng 2026:
\begin{itemize}
\item \textbf{Hệ sinh thái công nghệ} trưởng thành (Đà Nẵng IT Park, FPT Đà Nẵng,
Axon, Sun*, Enclave, nhiều outsourcing).
\item \textbf{Mạng lưới đại học} dày: ĐH Bách khoa, ĐH Kinh tế, ĐH Kiến trúc, ĐH
Duy Tân, ĐH FPT, ĐH Sư phạm, ĐH Ngoại ngữ.
\item \textbf{Ngành du lịch} lớn thứ 3 Việt Nam, với nhiều doanh nghiệp vừa đủ
lớn để cần tự động hóa.
\item \textbf{Chi phí vận hành} thấp hơn TP.HCM/Hà Nội 40--50\%, cho phép công ty
sống được với doanh thu thấp trong năm đầu.
\item \textbf{Không cạnh tranh trực tiếp} với hàng trăm startup AI ở Hà Nội/TP.HCM.
\end{itemize}

\subsection{Tại sao là lúc này (Why now)}

\begin{enumerate}
\item \textbf{Công nghệ đủ chín} — model inference rẻ, context window lớn, tool use ổn định.
\item \textbf{Khách hàng đủ chín} — đã qua giai đoạn ``AI là đồ chơi'', bắt đầu hỏi ``AI giúp tôi việc gì cụ thể?''.
\item \textbf{Chi phí vào nghề còn rẻ} — 2--3 năm nữa sẽ cần vốn lớn hơn nhiều để xây được vị thế.
\item \textbf{Cửa sổ vertical còn mở} — chưa có \textit{``Shopify cho admissions''}, \textit{``Zendesk cho nhà trường''}.
\item \textbf{Đủ kinh nghiệm cá nhân} — cả 3 người đều đã qua giai đoạn học việc, có đủ vốn kiến thức và quan hệ để bước vào vai trò founder.
\end{enumerate}

\subsection{Tại sao là 3 người chúng ta (Why us)}

\begin{itemize}
\item \textbf{Tổ hợp kỹ năng hiếm ở Việt Nam:}
\begin{itemize}
  \item \emph{Nghiên cứu \& định hướng} (Hảo — PhD Telecoms, giảng viên)
  \item \emph{Kỹ thuật AI \& phát triển} (Phú — học trò Hảo, đã có quá trình cộng tác)
  \item \emph{Kinh doanh \& thị trường} (Hoàng — doanh nhân có kinh nghiệm)
\end{itemize}
\item \textbf{Đã có lịch sử làm việc chung} (ít nhất giữa Hảo--Phú).
\item \textbf{Không trùng vai, không giẫm chân} — phân vai tự nhiên, không ép buộc.
\item \textbf{Có lợi thế network} ở cả phía nhà trường lẫn doanh nghiệp.
\item \textbf{Giai đoạn cuộc sống phù hợp} — đủ chín để nghiêm túc, chưa quá già để liều.
\end{itemize}

\clearpage

% =============================================================
% CHAPTER II — TẦM NHÌN & ĐÍCH ĐẾN
% =============================================================
\section{Tầm nhìn \& Đích đến}

\subsection{Tầm nhìn 10 năm}

\begin{quotebox}
\textit{Trong 10 năm, chúng tôi muốn trở thành công ty được các tổ chức tri thức
ở Đông Nam Á tin cậy nhất trong việc xây dựng hệ thống AI tự vận hành — nơi mà
khi một trường đại học, một hãng du lịch hay một tổ chức dịch vụ muốn ``agentify''
quy trình của họ, tên chúng tôi là tên đầu tiên họ nghĩ đến.}
\end{quotebox}

Đây là tầm nhìn ``công ty phục vụ'' chứ không phải ``công ty công nghệ khoe tech''.
Trọng tâm là \textit{tin cậy}, không phải \textit{ấn tượng}.

\subsection{Sứ mệnh}

\textbf{Giúp các tổ chức tri thức tại Việt Nam và Đông Nam Á tự động hóa những
việc lặp đi lặp lại một cách an toàn, minh bạch và có kiểm soát — để con người
của họ được giải phóng để làm việc con người nên làm.}

Chú ý: sứ mệnh không phải ``xây AI tốt nhất'', cũng không phải ``thay thế con người''.
Sứ mệnh là \textbf{giải phóng con người khỏi công việc máy nên làm}.

\subsection{Giá trị cốt lõi (5 giá trị)}

\begin{enumerate}
\item \textbf{Clarity over cleverness.} Rõ ràng hơn là khôn lỏi. Trong cả sản phẩm
lẫn giao tiếp. Một use case đơn giản chạy tốt $>$ một kiến trúc tuyệt vời nhưng
không ai hiểu.

\item \textbf{Operate, don't demo.} Làm cho chạy thật, không làm để trình diễn.
Mỗi dòng code phải có khách hàng thật đứng đằng sau nó.

\item \textbf{Trust is infrastructure.} Độ tin cậy là hạ tầng, không phải tính năng.
Evals, logs, audit trail, human-in-the-loop — không phải thứ làm sau.

\item \textbf{Long over loud.} Thà bền lâu hơn ồn ào. Chúng tôi chọn 10 năm phục vụ
trên 100 khách còn hơn 1 năm gọi vốn rầm rộ.

\item \textbf{Team before individual.} Không có ngôi sao. Ba người đồng sáng lập là
một đơn vị, không phải ba cá nhân đứng cạnh nhau.
\end{enumerate}

\subsection{Định vị thương hiệu}

\textbf{Định vị chính thức (outside-in):}\\
\textit{``AI Agent Systems for Knowledge-Intensive Workflows — built in Vietnam,
for Vietnam and Southeast Asia.''}

\textbf{Cách nói với khách hàng giáo dục:}\\
\textit{``Chúng tôi xây hệ thống AI giúp trường bạn tự xử lý hàng ngàn câu hỏi
và hàng trăm quy trình lặp lại mỗi tháng, mà vẫn giữ được chất lượng và
minh bạch.''}

\textbf{Cách nói với khách hàng du lịch / dịch vụ:}\\
\textit{``Chúng tôi xây hệ thống AI vận hành các quy trình tư vấn và dịch vụ
của bạn 24/7, có kiểm soát và có thể handoff cho người thật khi cần.''}

\textbf{Cái chúng tôi \underline{không} nói:}
\begin{itemize}
\item Không dùng từ ``framework'' khi nói với khách.
\item Không dùng ``multi-agent'' nếu không cần.
\item Không hứa ``thay thế con người''.
\item Không dùng buzzword (``revolutionary'', ``cutting-edge'', ``next-gen'') trong tài liệu chính thức.
\end{itemize}

\subsection{Mục tiêu định lượng: 1 năm / 3 năm / 5 năm}

\begin{tabularx}{\linewidth}{@{}p{1.8cm} X X X@{}}
\toprule
 & \textbf{Năm 1} (2026) & \textbf{Năm 3} (2028) & \textbf{Năm 5} (2030) \\
\midrule
\textbf{Khách hàng} & 2--3 design partner & 10--15 khách giáo dục & 50--100 khách (đa vertical) \\
\textbf{Doanh thu} & 500M--1B VND & 5--8B VND & 30--50B VND \\
\textbf{Nhân sự} & 3 founder + 1--2 hỗ trợ & 10--12 người & 30--50 người \\
\textbf{Sản phẩm} & 2 use case chạy thật & 5--6 solution packaged & Platform + subscription \\
\textbf{Vertical} & Giáo dục & Giáo dục + Du lịch thử & Đa vertical + địa bàn ĐNA \\
\textbf{Framework} & 6 khối cơ bản & Ổn định, có connector SDK & Được tái sử dụng 80\% dự án \\
\textbf{Thương hiệu} & Được biết trong giới giáo dục ĐN & Top-3 VN về agent systems & Top-3 ĐNA về agent systems \\
\bottomrule
\end{tabularx}

\begin{warnbox}[title=Cảnh báo]
Các con số trên là \textbf{ambition plausible}, không phải cam kết. Dùng để định
hướng quyết định, không phải để đánh giá thành công/thất bại. Sau 6 tháng
vận hành thực tế, các con số này \textbf{cần được hiệu chỉnh} dựa trên dữ liệu thật.
\end{warnbox}

\clearpage

% =============================================================
% CHAPTER III — CHIẾN LƯỢC SẢN PHẨM
% =============================================================
\section{Chiến lược sản phẩm}

\subsection{Triết lý ``Solution ngoài, Framework trong''}

Đây là nguyên tắc sản phẩm quan trọng nhất của công ty.

\textbf{Lớp ngoài (cái khách hàng mua):} Solution theo use case cụ thể, có ROI
đo được, có tên gọi dễ hiểu — ``Admissions Agent'', ``Student Services Agent'',
``Internal Knowledge Agent''.

\textbf{Lớp trong (tài sản của công ty):} Framework 6 khối, không bán rời, không
public, không đặt tên thương hiệu trong 12--18 tháng đầu.

\begin{center}
\begin{tikzpicture}[
  font=\sffamily,
  sol/.style={draw=ink, fill=accent!25, rounded corners=2pt,
              minimum width=36mm, minimum height=10mm, align=center, inner sep=2pt,
              font=\footnotesize\sffamily},
  solfut/.style={draw=ink!70, fill=accent!5, dashed, rounded corners=2pt,
              minimum width=36mm, minimum height=10mm, align=center, inner sep=2pt,
              font=\footnotesize\sffamily\itshape},
  fwk/.style={draw=ink, fill=mid!30, rounded corners=2pt,
              minimum width=26mm, minimum height=14mm, align=center, inner sep=2pt,
              font=\footnotesize\sffamily},
  arr/.style={-{Latex[length=3mm]}, thick, ink!75}
]
  % Layer A
  \node[sol] (a1) at (0,0) {Admissions\\Agent};
  \node[sol, right=3mm of a1] (a2) {Student Services\\Agent};
  \node[sol, right=3mm of a2] (a3) {Internal Knowledge\\Agent};
  \node[solfut, right=3mm of a3] (a4) {Service / Travel\\Agent (Year 2+)};

  \node[font=\footnotesize\sffamily\itshape, ink, above=1mm of a1.north west, anchor=south west]
    {LỚP A — Solutions bán ra thị trường (theo vertical + use case)};

  % Layer B
  \node[fwk, below=16mm of a1.south west, anchor=north west, text width=24mm] (b1)
    {\textbf{Knowledge}\\\scriptsize Ingest, metadata, versioning, citation};
  \node[fwk, right=2mm of b1, text width=24mm] (b2)
    {\textbf{Workflow}\\\scriptsize Intake, routing, state, approval};
  \node[fwk, right=2mm of b2, text width=24mm] (b3)
    {\textbf{Tools}\\\scriptsize Email, CRM, LMS, Drive, chat};
  \node[fwk, right=2mm of b3, text width=24mm] (b4)
    {\textbf{Reasoning}\\\scriptsize Single-agent, specialist, supervisor};
  \node[fwk, below=2mm of b2, text width=24mm] (b5)
    {\textbf{Evals}\\\scriptsize Logs, test set, confidence, feedback};
  \node[fwk, right=2mm of b5, text width=24mm] (b6)
    {\textbf{Governance}\\\scriptsize Access, audit, policy, approval};

  \node[font=\footnotesize\sffamily\itshape, ink, above=1mm of b1.north west, anchor=south west]
    {LỚP B — Core Framework (tài sản nội bộ, KHÔNG bán trực tiếp, KHÔNG đặt tên thương hiệu)};

  % Arrow B -> A
  \draw[arr] ($(b2.east)!0.5!(b3.west)$) ++(0,18mm) -- ++(0,4mm)
      node[right, font=\scriptsize\sffamily\itshape, ink!80]{\;đóng gói thành solution};
\end{tikzpicture}
\end{center}

\subsubsection*{Vì sao làm như vậy?}

\begin{itemize}
\item \textbf{Khách hàng không mua ``framework''} — họ mua kết quả. Nếu ta bán
framework, ta sẽ trở thành công ty tools cạnh tranh với LangChain, CrewAI,
AutoGen — cuộc chơi đó đã có quá nhiều người chơi.

\item \textbf{Framework nội bộ giúp ta nhanh hơn outsourcing agency} — mỗi khách
mới chỉ cần lắp 20--30\% phần custom, 70--80\% dùng lại. Đây là cách \textit{kinh
tế quy mô của chúng ta} được tạo ra.

\item \textbf{Framework là IP của công ty} — không phải bán ngay, mà là cái tạo ra
barrier to entry khi đối thủ muốn vào.
\end{itemize}

\subsection{Kiến trúc framework 6 khối}

Mỗi khối trả lời một câu hỏi cụ thể của agent system:

\begin{enumerate}
\item \textbf{Knowledge} — Agent biết gì? \textit{Ingest tài liệu, metadata,
versioning, citation, role-based access. Khách hàng cần trust rằng agent chỉ trả
lời từ dữ liệu chính thức của họ, không bịa.}

\item \textbf{Workflow} — Agent làm theo quy trình nào? \textit{Intake $\to$
routing $\to$ task state $\to$ escalation $\to$ approval. Đây là phần thay thế
\textbf{nghiệp vụ thủ công}, không phải phần AI.}

\item \textbf{Tools} — Agent dùng gì để làm việc? \textit{Email, form, CRM,
website, Drive, LMS/SIS, chat channels. Mỗi tool là một connector, phải có
standard interface.}

\item \textbf{Reasoning} — Agent suy nghĩ thế nào? \textit{Single-agent trước,
specialist agents sau, supervisor agent chỉ thêm khi thực sự cần. Đa số use case
giáo dục KHÔNG cần multi-agent.}

\item \textbf{Evaluation \& Observability} — Agent chạy có tốt không? \textit{
Log, test set, confidence threshold, human feedback, failure analysis. Đây là
thứ phân biệt sản phẩm thật với demo.}

\item \textbf{Governance} — Ai được làm gì, ai chịu trách nhiệm? \textit{Quyền
truy cập, audit trail, approval step, policy constraints. Đặc biệt quan trọng
với giáo dục và dịch vụ công.}
\end{enumerate}

\subsection{Chiến lược vertical: Giáo dục trước, Du lịch sau}

\subsubsection*{Vì sao giáo dục trước}

\begin{itemize}
\item Hảo có quyền truy cập và network tự nhiên trong ngành.
\item Quy trình tương đối chuẩn hóa (admissions, student services, academic knowledge).
\item Khách hàng có đủ budget và đủ nhu cầu, nhưng chưa có nhà cung cấp tốt.
\item Chu kỳ bán hàng không quá dài (3--6 tháng, không phải 12--18 tháng như tập đoàn).
\item ROI đo được (số câu hỏi tự động hóa, thời gian phản hồi, tỉ lệ escalation).
\end{itemize}

\subsubsection*{2 use case đầu trong giáo dục}

\textbf{Use case 1 — Admissions Agent.}
\begin{itemize}
\item \textit{Pain:} Phòng Tuyển sinh trả lời hàng ngàn câu hỏi lặp lại mùa cao điểm,
tư vấn ngành, hồ sơ, học phí, học bổng, chuyển ngành, chuyển trường.
\item \textit{Giải pháp:} Agent trả lời 24/7, truy vấn quy chế \& tài liệu chính
thức, escalate cho người thật khi cần, ghi log và đưa vào CRM.
\item \textit{ROI:} Giảm 60--80\% câu hỏi routine về phòng tuyển sinh, tăng tốc
độ phản hồi từ ngày xuống phút.
\end{itemize}

\textbf{Use case 2 — Internal Knowledge Agent.}
\begin{itemize}
\item \textit{Pain:} Giảng viên, cán bộ phải tìm kiếm quy định, văn bản, SOP,
thông báo rải rác ở nhiều nơi.
\item \textit{Giải pháp:} Agent trả lời nội bộ, có citation, truy cập theo vai
trò, có audit trail.
\item \textit{ROI:} Giảm 30--50\% thời gian tìm kiếm thông tin, tăng tuân thủ SOP.
\end{itemize}

\subsubsection*{Vì sao chưa đụng vertical khác}

\begin{itemize}
\item \textbf{Du lịch} — cần thêm knowledge về chính sách visa, tour, địa điểm;
chu kỳ bán hàng khác; KPI khác. \textit{Chờ đến Q3--Q4 năm đầu.}
\item \textbf{Y tế} — có pháp lý phức tạp. \textit{Không đụng trong 3 năm đầu.}
\item \textbf{Tài chính / ngân hàng} — compliance rất cao. \textit{Không đụng cho đến khi có người trong team hiểu sâu.}
\item \textbf{Chính phủ / khối nhà nước} — chu kỳ bán hàng quá dài, không phù hợp startup. \textit{Chỉ đụng khi có đối tác lớn.}
\end{itemize}

\subsection{Mô hình kinh doanh 3 giai đoạn}

\subsubsection*{Giai đoạn 1 (M1--M9): Design partner + pilot trả phí}

\begin{itemize}
\item Bán cái gì: Discovery workshop + workflow design + pilot + managed operations.
\item Giá: 50--200 triệu VND/pilot (tuỳ scope). Không miễn phí. Không dưới 20\% giá thật.
\item Số lượng: 2--3 design partner.
\item Mục tiêu: Có case study thật, có KPI thật, có reference thật.
\end{itemize}

\subsubsection*{Giai đoạn 2 (M10--M24): Solution packaging}

\begin{itemize}
\item Bán cái gì: 3 solution chuẩn (Admissions / Student Services / Internal Knowledge) + custom onboarding.
\item Giá: 200--500 triệu VND/trường/năm (license + support), hoặc mô hình hybrid.
\item Số lượng: 10--15 khách.
\item Mục tiêu: Doanh thu lặp lại, giảm thời gian triển khai xuống 4--6 tuần/khách.
\end{itemize}

\subsubsection*{Giai đoạn 3 (M25--M36+): Subscription / platformization}

\begin{itemize}
\item Bán cái gì: SaaS subscription cho các module chuẩn, + dịch vụ custom cho deployment lớn.
\item Giá: Theo tier (Starter / Pro / Enterprise), 100--800 triệu/năm.
\item Số lượng: 30--50+ khách.
\item Mục tiêu: Doanh thu bán buôn, biên lợi nhuận 40--60\%.
\end{itemize}

\begin{warnbox}[title=Không leo giai đoạn quá sớm]
Mỗi giai đoạn chỉ được bắt đầu khi giai đoạn trước có đủ \textbf{bằng chứng vận hành thật}:
\begin{itemize}
\item Không packaging khi chưa có 2--3 design partner chạy thật.
\item Không subscription khi chưa có 5+ khách dùng cùng solution ổn định.
\item Nếu ép giai đoạn, công ty sẽ bị biến thành agency mặc áo platform.
\end{itemize}
\end{warnbox}

\clearpage

% =============================================================
% CHAPTER IV — LỘ TRÌNH 36 THÁNG
% =============================================================
\section{Lộ trình 36 tháng}

\begin{center}
\begin{tikzpicture}[
  font=\sffamily, scale=1, every node/.style={font=\footnotesize\sffamily},
  phase/.style={rectangle, draw=ink, rounded corners=3pt, minimum height=12mm, align=center, inner sep=4pt},
  ms/.style={circle, draw=ink, fill=accent, minimum size=4mm, inner sep=0pt}
]
  % Timeline
  \draw[ink!60, very thick] (0,0) -- (15.5,0);
  \foreach \x/\m in {0/M1, 5/M12, 10/M24, 15/M36} {
    \draw[ink!60, thick] (\x,0.2) -- (\x,-0.2);
    \node[below=1mm] at (\x,-0.2) {\m};
  }

  % Year 1
  \node[phase, fill=hao!20, minimum width=48mm] at (2.5,1.4) {\textbf{Năm 1 — CHỨNG MINH}\\\tiny 1 vertical, 2 use case, 2--3 design partner};
  \node[ms] at (1,0) {}; \node[above=1mm, font=\tiny] at (1,0) {Prototype};
  \node[ms] at (2.5,0) {}; \node[above=1mm, font=\tiny] at (2.5,0) {1st paid};
  \node[ms] at (4,0) {}; \node[above=1mm, font=\tiny] at (4,0) {3 pilots};

  % Year 2
  \node[phase, fill=phu!20, minimum width=48mm] at (7.5,1.4) {\textbf{Năm 2 — NHÂN RỘNG}\\\tiny Packaging, 10--15 khách, bắt đầu thử du lịch};
  \node[ms] at (6,0) {}; \node[above=1mm, font=\tiny] at (6,0) {5 khách};
  \node[ms] at (8,0) {}; \node[above=1mm, font=\tiny] at (8,0) {Hire \#4};
  \node[ms] at (9.5,0) {}; \node[above=1mm, font=\tiny] at (9.5,0) {10 khách};

  % Year 3
  \node[phase, fill=hoang!20, minimum width=48mm] at (12.5,1.4) {\textbf{Năm 3 — PLATFORM}\\\tiny Subscription, đa vertical, ĐNA};
  \node[ms] at (11,0) {}; \node[above=1mm, font=\tiny] at (11,0) {Sub v1};
  \node[ms] at (13,0) {}; \node[above=1mm, font=\tiny] at (13,0) {25 khách};
  \node[ms] at (15,0) {}; \node[above=1mm, font=\tiny] at (15,0) {50 khách};
\end{tikzpicture}
\end{center}

\subsection{Năm 1 (2026) — CHỨNG MINH}

\textbf{Câu hỏi lớn:} Có khách hàng nào trả tiền thật cho chúng ta không, và
chúng ta có vận hành được không?

\begin{tabularx}{\linewidth}{@{}p{1.2cm} X X@{}}
\toprule
\textbf{Quý} & \textbf{Mục tiêu} & \textbf{Deliverable} \\
\midrule
Q1 & Pháp lý, framework v0, prototype & Cty thành lập; prototype Admissions Agent chạy trên 1 trường thật (pilot nội bộ); 10 discovery call. \\
Q2 & Design partner đầu tiên & 1 khách ký hợp đồng pilot trả phí; framework 6 khối v0.5; case study v1. \\
Q3 & Mở rộng design partner & 2--3 khách chạy pilot; framework v1.0 ổn định; bắt đầu viết ``playbook delivery'' chuẩn. \\
Q4 & Hoàn thiện solution \#1 \& \#2 & Admissions Agent \& Internal Knowledge Agent đã được packaging đủ để bán cho khách \#4. \\
\bottomrule
\end{tabularx}

\textbf{KPI năm 1:}
\begin{itemize}
\item $\geq$ 2 khách trả phí (tối thiểu 50M/khách).
\item $\geq$ 10 discovery call có ghi chép.
\item Framework 6 khối: có v1.0.
\item Doanh thu: 500M--1B VND.
\item Không chết tiền mặt (burn rate không vượt quá thu nhập riêng của 3 founder đóng góp).
\end{itemize}

\subsection{Năm 2 (2027) — NHÂN RỘNG}

\textbf{Câu hỏi lớn:} Chúng ta có thể làm được 10--15 khách mà không bị quá tải?

\begin{itemize}
\item \textbf{Q1--Q2:} Tuyển 1--2 người (kỹ sư \#2, customer success). Packaging
3 solution chuẩn. Thời gian triển khai mỗi khách mới xuống 6--8 tuần.
\item \textbf{Q3--Q4:} Bắt đầu thử vertical 2 (du lịch / dịch vụ) với 1 khách pilot.
Không đầu tư lớn. Dùng framework sẵn có để kiểm chứng tính tổng quát.
\end{itemize}

\textbf{KPI năm 2:}
\begin{itemize}
\item 10--15 khách giáo dục trả phí.
\item Doanh thu 3--5B VND.
\item Nhân sự: 8--12 người.
\item Biên lợi nhuận gộp: 30--40\%.
\item Framework đã được dùng trên 100\% dự án mới.
\end{itemize}

\subsection{Năm 3 (2028) — PLATFORM HÓA}

\textbf{Câu hỏi lớn:} Chúng ta có thể chuyển từ dịch vụ sang subscription mà không mất khách?

\begin{itemize}
\item \textbf{Q1--Q2:} Ra mắt subscription v1 cho Admissions Agent \& Internal Knowledge Agent. Chuyển 30--50\% khách hiện tại sang subscription.
\item \textbf{Q3--Q4:} Mở rộng sang Đông Nam Á (1 khách Philippines / Indonesia / Thailand). Thử nghiệm reseller model ở ĐNA.
\end{itemize}

\textbf{KPI năm 3:}
\begin{itemize}
\item 30--50 khách tổng.
\item Doanh thu 15--25B VND.
\item Nhân sự: 20--30 người.
\item $\geq$ 40\% doanh thu từ subscription.
\item 1 khách ngoài Việt Nam.
\end{itemize}

\subsection{Decision gates — 4 mốc quyết định lớn}

\begin{tabularx}{\linewidth}{@{}p{1.5cm} p{3cm} X X@{}}
\toprule
\textbf{Mốc} & \textbf{Thời điểm} & \textbf{Câu hỏi} & \textbf{Nếu KHÔNG đạt thì làm gì} \\
\midrule
\textbf{Gate 1} & M6 & Có 1 khách trả phí chưa? & Giảm quy mô, Hảo \& Hoàng giữ công việc cũ thêm 6 tháng, Phú về part-time. \\
\textbf{Gate 2} & M12 & Có 3 khách và framework v1.0 chưa? & Dừng thử thêm vertical. Tập trung hoàn thiện solution \#1 trong 6 tháng nữa. \\
\textbf{Gate 3} & M24 & Có 10+ khách và biên lợi nhuận 30\%+ chưa? & Dừng mở vertical mới. Không tuyển thêm. Tập trung profitability. \\
\textbf{Gate 4} & M36 & Có subscription chạy và 40\% khách chuyển đổi chưa? & Lùi kế hoạch ĐNA. Tiếp tục là công ty dịch vụ \& solution tại VN. \\
\bottomrule
\end{tabularx}

\begin{warnbox}[title=Nguyên tắc decision gate]
Tại mỗi gate, 3 người ngồi lại, nhìn dữ liệu, và quyết định dựa trên thực tế.
\textbf{Không được bỏ qua gate.} Không được tự động leo giai đoạn tiếp theo vì
\textit{``cảm thấy ổn''}.
\end{warnbox}

\clearpage

% =============================================================
% CHAPTER V — CON NGƯỜI
% =============================================================
\section{Con người}

Đây là chương quan trọng nhất của playbook. Chiến lược có thể sai rồi sửa.
Sản phẩm có thể đi sai hướng rồi điều chỉnh. Nhưng nếu \textbf{3 người không
đồng thuận về vai trò, quyền, trách nhiệm, và cách cư xử với nhau}, công ty
sẽ vỡ trước khi sản phẩm kịp chín.

\subsection{Triết lý đội ngũ: ``Ba con người, một đơn vị''}

\begin{quotebox}
Chúng tôi không phải 3 cá nhân hợp tác. Chúng tôi là \textbf{một đơn vị 3 đầu}.
Quyết định của công ty không phải ``đa số 2/3 thắng'', mà là \textbf{sự đồng thuận
có nguyên tắc}, trong đó mỗi người có vùng quyết định chính của riêng mình, và
có nghĩa vụ tôn trọng vùng của người khác.
\end{quotebox}

Điều này có nghĩa:
\begin{itemize}
\item Không có \textit{``CEO độc tài''} — dù Hảo là người ký giấy tờ chính thức.
\item Không có \textit{``kỹ sư làm theo lệnh''} — dù Phú nhận trách nhiệm delivery.
\item Không có \textit{``sales nói gì cũng đúng''} — dù Hoàng tiếp xúc khách hàng.
\end{itemize}

\subsection{Đỗ Phúc Hảo — CEO / Strategy \& Direction}

\subsubsection*{Background}
PhD Telecoms (Russia), hiện giảng dạy tại ĐH Kiến trúc Đà Nẵng. Có nền tảng
nghiên cứu sâu, kinh nghiệm giảng dạy, và mạng lưới học thuật rộng. Đã và đang
tham gia các đề án nghiên cứu, viết báo khoa học, cũng như các dự án phần mềm
giáo dục (KĐCLGD, CTĐT redesign).

\subsubsection*{Ngũ hành \& tính cách}
Canh Ngọ --- Lộ Bàng Thổ --- giờ Dậu. Đất ven đường --- điểm đi lại, nền móng,
kết nối các phía. Thổ sinh Kim (hỗ trợ Phú tự nhiên). Vai trò trục --- giữ cho
công ty không lệch hướng.

\subsubsection*{Trách nhiệm chính}
\begin{enumerate}
\item \textbf{Chọn vertical \& chọn use case.} Quyết định cuối cùng về công ty
đi vào ngành nào, làm use case nào.
\item \textbf{Giữ roadmap sản phẩm.} Nói KHÔNG với những thứ không nằm trong
roadmap, kể cả khi Hoàng có khách muốn mua.
\item \textbf{Gỡ xung đột kỹ thuật--sales.} Khi Phú và Hoàng không đồng ý, Hảo
là người chốt.
\item \textbf{Đại diện công ty trong giới học thuật.} Paper, hội thảo, đề án
nghiên cứu chung với trường.
\item \textbf{Hiring \& văn hóa.} Phỏng vấn lần cuối mọi vị trí, giữ chuẩn
chất lượng người.
\item \textbf{Gây vốn (nếu cần).} Đàm phán với nhà đầu tư, đối tác chiến lược.
\end{enumerate}

\subsubsection*{Quyền quyết định (decision rights)}
\begin{itemize}
\item \textbf{Solo:} Hiring cuối cùng; đi hay không đi hội thảo/sự kiện; hợp tác với trường/viện.
\item \textbf{Shared với Phú:} Roadmap kỹ thuật (3 tháng, 6 tháng).
\item \textbf{Shared với Hoàng:} Chọn ngành tiếp theo; chính sách giá.
\item \textbf{Cần 3 chữ ký:} Ngân sách $>$100 triệu; đổi định vị công ty; thêm đồng sáng lập.
\end{itemize}

\subsubsection*{KPI (năm 1)}
\begin{itemize}
\item Vertical đã chọn có $\geq$ 2 design partner trả phí.
\item $\geq$ 60\% code framework được dùng lại qua $\geq$ 2 dự án.
\item $\geq$ 1 bài viết / talk công khai xây thương hiệu học thuật công ty.
\item 100\% decision gate được đánh giá đúng hạn, có số liệu.
\end{itemize}

\subsubsection*{Những thứ Hảo KHÔNG làm}
\begin{itemize}
\item Không viết code production (có thể prototype, không commit vào main repo).
\item Không trực tiếp đi sales lạnh (có thể giới thiệu warm intro).
\item Không làm delivery PM (có thể review, không own timeline).
\item Không quyết định kiến trúc kỹ thuật (Phú quyết, Hảo có quyền chất vấn).
\end{itemize}

\subsubsection*{Growth path 3 năm}
Từ CEO-founder $\to$ CEO có CTO \& COO dưới quyền. Vai trò dần chuyển về
strategic partnership, đại diện công ty ở ngoài, và hướng dẫn R\&D dài hạn.

\subsection{Lê Hữu Phú — Platform Lead / Head of Engineering}

\subsubsection*{Background}
Học trò của Hảo, đã từng làm việc cùng Hảo trong các dự án. Có nền tảng kỹ
thuật, được đào tạo dưới sự dẫn dắt của Hảo --- đây là \textbf{lợi thế rất lớn}
của công ty: kỹ sư chính đã có tiếng nói chung với người định hướng.

\subsubsection*{Ngũ hành \& tính cách}
Canh Thìn --- Bạch Lạp Kim --- giờ Ngọ. ``Kim sáp'' --- linh hoạt, có khả năng
định hình, được Hảo (Thổ) hỗ trợ tự nhiên. Hỏa của giờ Ngọ tạo độ bứt phá,
nhưng cũng cần tránh ``tiêu kim'' --- tức là làm nhiều thứ cùng lúc quá mức.

\subsubsection*{Trách nhiệm chính}
\begin{enumerate}
\item \textbf{Xây core framework.} Đây là di sản lớn nhất của Phú trong công ty.
6 khối framework phải ra hình sau năm 1.
\item \textbf{Giữ chất lượng kiến trúc.} Chống lại xu hướng ``patch nhanh cho
khách này'' mà làm hỏng framework.
\item \textbf{Chuẩn hóa quy trình delivery.} Để khi có kỹ sư \#2, \#3, họ onboard
được trong 2--3 tuần.
\item \textbf{Giữ tốc độ MVP.} Mỗi use case mới: prototype trong 2 tuần, pilot
trong 8 tuần.
\item \textbf{Chi phí vận hành hệ thống.} Giữ cost/khách ở mức bền vững.
\item \textbf{Evals \& observability.} Đây không phải ``tính năng phụ'' --- đây
là phần khiến khách hàng tin công ty.
\end{enumerate}

\subsubsection*{Quyền quyết định}
\begin{itemize}
\item \textbf{Solo:} Tech stack; coding standards; kiến trúc module; tool \& CI/CD.
\item \textbf{Shared với Hảo:} Roadmap kỹ thuật; hire kỹ sư.
\item \textbf{Shared với Hoàng:} Thời gian delivery của mỗi dự án (không tự cam kết với khách).
\item \textbf{Cần 3 chữ ký:} Đổi foundation model chính; rewrite kiến trúc lớn; tích hợp với bên thứ 3 có pháp lý.
\end{itemize}

\subsubsection*{KPI (năm 1)}
\begin{itemize}
\item Framework 6 khối có v1.0 sau 9 tháng.
\item Prototype 1 use case trong 2 tuần.
\item Pilot 1 use case trong 8 tuần.
\item Uptime các pilot $\geq$ 99\%.
\item Zero data leak / security incident.
\item $\geq$ 60\% module tái sử dụng giữa dự án \#1 và dự án \#3.
\end{itemize}

\subsubsection*{Những thứ Phú KHÔNG làm}
\begin{itemize}
\item Không trực tiếp bán hàng với khách (có thể tham gia kỹ thuật meeting).
\item Không hứa tính năng với khách khi Hoàng không có mặt.
\item Không mở rộng scope dự án mà không có 3 chữ ký.
\item Không từ chối use case với lý do kỹ thuật mà không cho lựa chọn thay thế.
\end{itemize}

\subsubsection*{Growth path 3 năm}
Từ sole engineer $\to$ Head of Engineering của team 6--10 người $\to$ CTO với
nhóm research, nhóm platform, nhóm delivery. Ngày Phú có kỹ sư \#2 là ngày
Phú phải học kỹ năng mới quan trọng nhất: \textbf{mentoring và code review}.

\begin{warnbox}[title=Rủi ro lớn nhất của Phú]
Làm một mình quá lâu. Nếu đến M6 mà chưa có kỹ sư \#2 giúp đỡ ít nhất part-time,
Phú sẽ \textbf{burnout}. Đây là rủi ro lớn nhất của công ty, không phải rủi ro kỹ
thuật hay rủi ro khách hàng. \textbf{Hảo và Hoàng có trách nhiệm theo dõi energy
của Phú và không để việc này xảy ra.}
\end{warnbox}

\subsection{Nguyễn Ngọc Hoàng — GTM Lead / Head of Business}

\subsubsection*{Background}
Doanh nhân có nhiều năm kinh nghiệm. Hiểu quy trình bán hàng B2B, có mạng
lưới doanh nghiệp tại Đà Nẵng và miền Trung. Khác với Hảo và Phú
(dân kỹ thuật/học thuật), Hoàng mang đến \textbf{góc nhìn thị trường thực dụng}
--- cái mà công ty nào thuần kỹ thuật cũng thiếu.

\subsubsection*{Ngũ hành \& tính cách}
Kỷ Tỵ --- Đại Lâm Mộc --- giờ Dần/Mão. Rừng lớn, cây cổ thụ --- bao quát, có tán
rộng, giỏi xây mạng lưới. Mộc vượng khi gặp Mộc giờ. Tính cách: cần không gian
để thở, không hợp với micromanage chi tiết.

\subsubsection*{Trách nhiệm chính}
\begin{enumerate}
\item \textbf{Tìm pain point thật.} Không phải ``khách hàng bảo họ muốn gì'' --- mà
là ``khách hàng thật sự đang khổ vì gì''.
\item \textbf{Xác định ICP (Ideal Customer Profile).} Ai là trường sẵn sàng trả
tiền? Bao nhiêu sinh viên? Ngân sách bao nhiêu? Ai là người ra quyết định?
\item \textbf{Chốt design partner.} Chuyển từ ``quan tâm'' sang ``trả tiền''.
\item \textbf{Dẫn discovery workshop.} Hoàng đưa Phú vào cuộc, không làm một mình.
\item \textbf{Pricing \& đàm phán hợp đồng.} Không giảm giá dưới 20\% base rate mà không có 3 chữ ký.
\item \textbf{Customer success sau ký hợp đồng.} Giữ khách, upsell khi phù hợp.
\item \textbf{Insight thị trường.} Đều đặn cập nhật 3 người về xu hướng ngành.
\end{enumerate}

\subsubsection*{Quyền quyết định}
\begin{itemize}
\item \textbf{Solo:} Kỹ thuật discovery; chiến thuật tiếp cận khách; scheduling sales.
\item \textbf{Shared với Hảo:} Chọn ngành tiếp theo; chính sách giá chung.
\item \textbf{Shared với Phú:} Scope cụ thể của mỗi hợp đồng (Phú xác nhận làm được).
\item \textbf{Cần 3 chữ ký:} Ký hợp đồng $>$300 triệu; thử vertical mới; đối tác reseller.
\end{itemize}

\subsubsection*{KPI (năm 1)}
\begin{itemize}
\item $\geq$ 10 discovery call có ghi chép đầy đủ.
\item $\geq$ 2 khách ký pilot trả phí.
\item ICP rõ ràng sau M4.
\item Pipeline $\geq$ 10 qualified leads sau M9.
\item Tỷ lệ use case $\to$ proposal $\geq$ 30\%.
\end{itemize}

\subsubsection*{Những thứ Hoàng KHÔNG làm}
\begin{itemize}
\item Không hứa tính năng chưa có 3 chữ ký.
\item Không giảm giá dưới 20\% base rate không có Hảo đồng ý.
\item Không nhận dự án ngoài vertical đã chọn (không có 3 chữ ký).
\item Không tự quyết scope kỹ thuật (phải có Phú xác nhận).
\item Không tiếp cận khách mà không có material chuẩn (không đi ``tay không'').
\end{itemize}

\subsubsection*{Growth path 3 năm}
Từ sole business lead $\to$ Head of Business với team sales + customer success
(4--6 người) $\to$ COO / Chief Revenue Officer. Kỹ năng cần học: từ ``bán hàng''
sang ``xây team bán hàng''.

\subsection{Ngũ hành: sự bổ sung có chủ ý}

\begin{center}
\begin{tikzpicture}[
  font=\sffamily,
  role/.style={circle, draw=ink, thick, minimum size=22mm, align=center, font=\footnotesize\sffamily\bfseries},
  arrgood/.style={-{Latex[length=3mm]}, thick},
  arrwarn/.style={-{Latex[length=3mm]}, thick, densely dashed}
]
  \node[role, fill=hao!30] (h) at (0,0) {Hảo\\Thổ\\\scriptsize Strategy};
  \node[role, fill=phu!40] (p) at (5,0) {Phú\\Kim\\\scriptsize Platform};
  \node[role, fill=hoang!30] (ho) at (10,0) {Hoàng\\Mộc\\\scriptsize GTM};

  \draw[arrgood, hao]   (h)  to[bend left=18]  node[above, font=\scriptsize, hao]{Thổ sinh Kim --- nuôi dưỡng tự nhiên} (p);
  \draw[arrwarn, warm]  (p)  to[bend left=18]  node[above, font=\scriptsize, warm]{Kim khắc Mộc --- cần Hảo làm trọng tài} (ho);
  \draw[arrgood, hoang] (ho) to[bend left=30]  node[below, font=\scriptsize, hoang]{Mộc hút Thổ --- Hoàng leverage network của Hảo} (h);
\end{tikzpicture}
\end{center}

\textbf{Diễn giải cho nội bộ:}

\begin{itemize}
\item \textbf{Quan hệ thuận --- Hảo $\to$ Phú.} Đây là trục sản phẩm tự nhiên:
Hảo định hướng, Phú kết tinh thành sản phẩm. Không cần cơ chế đặc biệt, chỉ cần
tôn trọng lẫn nhau.

\item \textbf{Quan hệ ngang --- Hoàng $\to$ Hảo.} Hoàng tận dụng network học
thuật và uy tín của Hảo để mở cửa khách hàng đầu. Đổi lại, Hoàng bảo vệ Hảo
khỏi các yêu cầu không hợp lý từ khách hàng.

\item \textbf{Điểm cần chú ý --- Phú vs Hoàng.} Kỹ thuật và sales luôn có mâu
thuẫn tiềm năng: sales muốn cam kết nhiều, kỹ thuật muốn làm ít--chắc. Đây là
lý do \textbf{cơ chế 3 chữ ký} tồn tại. Hảo là người trung gian bắt buộc.
\end{itemize}

\subsection{Cơ chế ra quyết định}

\subsubsection*{3 tầng quyết định}

\begin{tabularx}{\linewidth}{@{}p{2.2cm} X X@{}}
\toprule
\textbf{Tầng} & \textbf{Ai quyết} & \textbf{Ví dụ} \\
\midrule
\textbf{Vận hành hằng ngày} & Mỗi người trong vùng của mình & Phú chọn library; Hoàng chọn thời gian gặp khách; Hảo duyệt bài đăng LinkedIn. \\
\textbf{Quyết định vừa} & 2 người liên quan + thông báo người thứ 3 & Pricing 1 hợp đồng (Hoàng + Hảo); scope dự án (Phú + Hoàng); roadmap quý (Hảo + Phú). \\
\textbf{Quyết định lớn --- ``3 chữ ký''} & Phải có đồng thuận 3 người, bằng văn bản & Ngân sách $>$100M; đổi vertical; đổi định vị; hire/fire; reshape cổ phần; vay nợ; thuê nhà. \\
\bottomrule
\end{tabularx}

\subsubsection*{Khi 3 người không đồng ý}

\begin{enumerate}
\item \textbf{Dừng lại 48h.} Không quyết định nóng.
\item \textbf{Viết quan điểm.} Mỗi người viết 1 trang: tôi đề xuất gì, vì sao, và nếu sai thì hậu quả gì.
\item \textbf{Ngồi lại.} Đọc của nhau trước. Nói sau.
\item \textbf{Nếu vẫn không đồng ý:} Thử pilot với scope nhỏ, timebox 4--8 tuần, đánh giá bằng dữ liệu.
\item \textbf{Nếu không thể timebox:} Xếp lại quyết định sau 1--2 tháng, quay lại khi có thêm thông tin.
\end{enumerate}

\begin{warnbox}[title=Luật không thành văn]
Không ai được dùng ``tôi là cổ đông lớn hơn'' hoặc ``tôi là CEO'' làm argument
cuối cùng. Nếu phải dùng đến chức danh để thắng, \textbf{quyết định đó đã sai rồi}.
\end{warnbox}

\subsection{Cổ phần, vốn, và vesting (đề xuất để thảo luận)}

Đây là đề xuất \textbf{khung để đối thoại}, không phải quyết định. 3 người cần
ngồi lại riêng một buổi để chốt.

\subsubsection*{Đề xuất cổ phần khởi đầu}

\textbf{Option A --- Chia đều 33 / 33 / 33 (khuyến nghị).}
\begin{itemize}
\item Ưu: thông điệp rõ ``một đơn vị'', không ai làm chủ ai.
\item Nhược: không phản ánh khác biệt về đóng góp giá trị (role criticality, kinh nghiệm, vốn).
\item \textbf{Nên chọn nếu:} cả 3 người đều full-time, không ai góp vốn lớn hơn hẳn người khác, và tin nhau cơ bản.
\end{itemize}

\textbf{Option B --- Chia theo trọng số.}
\begin{itemize}
\item Ví dụ 40 / 30 / 30 nếu Hảo là người dẫn dắt định hướng và đóng vai trò thu hút 2 người còn lại.
\item Ví dụ 35 / 35 / 30 nếu Hoàng tham gia part-time trong 6 tháng đầu.
\item \textbf{Nên chọn nếu:} có khác biệt thật về cam kết thời gian hoặc vốn góp.
\end{itemize}

\subsubsection*{Vesting schedule (khuyến nghị)}

\begin{itemize}
\item \textbf{4 năm vesting, 1 năm cliff.} Nghĩa là: nếu bất kỳ ai bỏ cuộc trong
năm đầu, họ \textbf{không được giữ cổ phần}. Sau năm 1, cổ phần được vest dần
hàng quý.
\item Đây không phải lời dọa --- đây là \textbf{bảo hiểm cho 2 người ở lại}. Nếu
1 người bỏ sau M3 với 33\% cổ phần, công ty sẽ không thể gọi vốn hay mời
người mới.
\end{itemize}

\subsubsection*{Vốn góp ban đầu (đề xuất)}

\begin{itemize}
\item Mỗi người góp một khoản tiền mặt nhỏ (ví dụ 50--100 triệu VND) để tạo runway
6--9 tháng cho chi phí cơ bản (thuê văn phòng nhỏ, cloud, công cụ, pháp lý).
\item Nếu một người góp nhiều hơn đáng kể, phần chênh nên được xử lý như
\textbf{convertible note}, không trộn vào cổ phần.
\item Không vay nợ cá nhân để bơm vào công ty trong 6 tháng đầu --- rủi ro quá lớn.
\end{itemize}

\subsubsection*{Lương trong năm đầu}

\begin{itemize}
\item \textbf{M1--M6:} Không có lương, hoặc lương tối thiểu (4--6 triệu/tháng) để
đóng bảo hiểm. Tất cả dồn vào runway.
\item \textbf{M7--M12:} Khi có doanh thu thật, bắt đầu trả lương cơ bản
(10--15 triệu/tháng).
\item \textbf{Năm 2+:} Lương thị trường. Ưu tiên trả Phú ngang hoặc cao hơn
thị trường kỹ sư AI (tránh mất Phú).
\end{itemize}

\subsection{Văn hóa \& cam kết}

\subsubsection*{Cam kết thời gian}

\begin{itemize}
\item Cả 3 người cam kết \textbf{tối thiểu 3 năm} đi cùng công ty.
\item Nếu 1 người muốn rút sớm hơn, phải thông báo trước \textbf{3 tháng}, và hỗ trợ handoff.
\item Tuyệt đối không có ``đang làm founder ở đây, vừa mở công ty khác''.
\end{itemize}

\subsubsection*{Cam kết minh bạch}

\begin{itemize}
\item \textbf{Tài chính công ty}: cả 3 có quyền xem mọi giao dịch, mọi hợp đồng.
\item \textbf{Lịch làm việc}: có lịch chung, ai làm gì tuần này đều có thể thấy.
\item \textbf{Quyết định}: ghi chép trong 1 Google Doc / Notion chung. Không có
quyết định ``trong đầu''.
\end{itemize}

\subsubsection*{Cam kết cá nhân --- ``5 điều không làm''}

\begin{enumerate}
\item Không nói xấu đồng sáng lập sau lưng.
\item Không hứa với khách hàng điều 2 người còn lại không biết.
\item Không tuyển người thân trong năm đầu (tránh nhạy cảm).
\item Không đầu tư cá nhân vào khách hàng hoặc đối thủ của công ty.
\item Không ra quyết định quan trọng khi đang tức giận, mệt mỏi, hoặc say.
\end{enumerate}

\subsubsection*{Nhịp họp}

\begin{tabularx}{\linewidth}{@{}p{3cm} p{2.5cm} X@{}}
\toprule
\textbf{Họp} & \textbf{Tần suất} & \textbf{Mục đích} \\
\midrule
\textbf{Daily standup} & 15 phút, 9h sáng & Hôm nay làm gì, kẹt gì. \\
\textbf{Weekly review} & 90 phút, thứ 6 & Metrics tuần, pipeline, kỹ thuật. \\
\textbf{Monthly strategy} & Nửa ngày, cuối tháng & Roadmap, pricing, hiring. \\
\textbf{Quarterly offsite} & 1 ngày + 1 bữa tối & Nhìn lại quý, nhìn tới quý tới. \\
\textbf{Annual retreat} & 2--3 ngày, ngoài DN & Nhìn lại năm, hiệu chỉnh 3 năm. \\
\textbf{1:1 cặp đôi} & 30 phút, 2 tuần/lần & Mỗi cặp 2 người nói chuyện riêng. \\
\bottomrule
\end{tabularx}

\clearpage

% =============================================================
% CHAPTER VI — RỦI RO & PHƯƠNG ÁN
% =============================================================
\section{Rủi ro \& Phương án}

\subsection{10 rủi ro lớn nhất}

\begin{tabularx}{\linewidth}{@{}p{0.5cm} p{3.3cm} X X@{}}
\toprule
\textbf{\#} & \textbf{Rủi ro} & \textbf{Xác suất} & \textbf{Phương án giảm thiểu} \\
\midrule
1 & Không tìm được design partner trả phí trong 6 tháng & Trung bình & Giảm giá pilot, mở rộng tới 5 trường khả thi, dùng network Hảo nâng độ tin cậy. \\
2 & Phú burnout do làm một mình quá lâu & \textbf{Cao} & Tuyển kỹ sư \#2 part-time sau M6; Hảo theo dõi energy; cấm làm cuối tuần thường xuyên. \\
3 & Hoàng không quen bán AI (khác với business cũ) & Trung bình & Hảo đi cùng 10 discovery đầu; có playbook sales chuẩn; có tài liệu kỹ thuật dễ hiểu. \\
4 & Mâu thuẫn Phú--Hoàng về scope & \textbf{Cao} & Cơ chế 3 chữ ký; Hảo làm trọng tài; timebox pilot thay vì tranh luận. \\
5 & Big tech (FPT / VNG / OpenAI / Anthropic) tung giải pháp vertical & Trung bình & Tập trung ngách Việt hoá \& hiểu sâu nghiệp vụ; xây moat bằng design partner trung thành. \\
6 & Model cost tăng đột ngột & Thấp--Trung & Xây cost tracking từ M1; sẵn sàng chuyển giữa 2--3 model; self-host model nhỏ cho task đơn giản. \\
7 & Khách thứ 3, 4 không cho tái sử dụng code vì custom quá nhiều & Trung bình & Hảo quyết định ngay từ pitch: công ty không nhận dự án không reuse được $\geq$ 60\%. \\
8 & Cạn tiền mặt trước khi có doanh thu lặp lại & \textbf{Cao} & Gate 1 ở M6; mỗi người giữ nguồn thu phụ; không thuê văn phòng đắt trong năm đầu. \\
9 & Một trong 3 người phải về việc cũ (ngoài ý muốn) & Thấp & Vesting 4 năm + 1 năm cliff; bảo hiểm nhân thọ cơ bản cho mỗi founder. \\
10 & Khách hàng mất niềm tin sau 1 lỗi công khai & Thấp--Trung & Evals \& audit trail ngay từ MVP; quy trình handle incident rõ ràng; không bao giờ hứa 100\%. \\
\bottomrule
\end{tabularx}

\subsection{Red lines --- những việc không bao giờ làm}

\begin{enumerate}
\item \textbf{Không nhận dự án không thể tái sử dụng $\geq$ 60\% framework.}
Ngay cả khi khách trả nhiều. Đó là outsourcing trá hình.

\item \textbf{Không bán sản phẩm chưa chạy thật.} Không có demo ``fake it till
you make it''. Demo là MVP thật, hoặc không demo.

\item \textbf{Không làm dự án có yêu cầu bỏ qua evals / audit / governance.}
Không phải vì đạo đức cao siêu --- vì nó sẽ gây incident sau này.

\item \textbf{Không hứa với khách điều 2 người còn lại chưa biết.} Mọi cam kết
phải qua 3 chữ ký.

\item \textbf{Không pha loãng cổ phần founder dưới 51\% trong 24 tháng đầu.}
Không gọi vốn lớn nếu chưa đủ traction.

\item \textbf{Không dùng tiền khách hàng cho việc ngoài dự án của khách đó.}
Không ``lấy tiền dự án A trả lương làm dự án B''.

\item \textbf{Không cover-up lỗi.} Với khách, với team, với nhau.
\end{enumerate}

\subsection{Go / No-go criteria --- khi nào thì DỪNG}

\begin{warnbox}[title=Đây là phần khó nhất]
Đa số startup không chết vì không biết bắt đầu, mà vì không biết dừng. Dưới đây
là các tín hiệu mà nếu xuất hiện, 3 người phải ngồi lại nghiêm túc:
\end{warnbox}

\begin{itemize}
\item \textbf{M6, chưa có 1 khách trả phí:} Không fail, nhưng cần phân tích nghiêm
túc. Có thể kéo dài thêm 3 tháng.

\item \textbf{M9, chưa có 1 khách trả phí:} Cảnh báo đỏ. Phải quyết định:
tiếp tục 3 tháng nữa với scope hẹp hơn, hay downgrade công ty thành agency.

\item \textbf{M12, vẫn chưa có khách trả phí:} \textbf{Dừng lại.} Không phải
thất bại cá nhân --- là thị trường đã nói. Giữ framework làm tài sản, giữ quan
hệ 3 người, nhưng tạm dừng mô hình công ty này.

\item \textbf{Bất kỳ lúc nào: một founder bị stress nghiêm trọng, ảnh hưởng sức
khỏe hoặc gia đình:} Dừng mọi việc khác, họp 3 người, ưu tiên con người trước công ty.

\item \textbf{Bất kỳ lúc nào: đã có 3 lần vi phạm 3-chữ-ký:} Nghĩa là cơ chế đã
gãy. Họp khẩn cấp. Không được ``thôi bỏ qua lần này''.
\end{itemize}

\clearpage

% =============================================================
% CHAPTER VII — 90 NGÀY ĐẦU & CAM KẾT
% =============================================================
\section{90 ngày đầu \& Cam kết}

\subsection{Kế hoạch 90 ngày (M1--M3)}

\subsubsection*{Tháng 1 --- Khởi động}

\textbf{Tuần 1--2:}
\begin{itemize}
\item Đọc, bình luận, hiệu chỉnh playbook này. Chốt phiên bản v1.0.
\item Ngồi 1 ngày riêng để chốt cổ phần, vesting, vốn góp, lương.
\item Viết biên bản đồng thuận 3 người (có chữ ký).
\end{itemize}

\textbf{Tuần 3--4:}
\begin{itemize}
\item Pháp lý: đăng ký công ty TNHH (2 hoặc 3 thành viên), địa chỉ Đà Nẵng.
\item Mở tài khoản ngân hàng doanh nghiệp.
\item Thuê văn phòng nhỏ (co-working cũng được, 3--5 chỗ ngồi).
\item Domain, email, công cụ làm việc (Notion, GitHub, Slack).
\end{itemize}

\subsubsection*{Tháng 2 --- Xây dựng \& Discovery}

\textbf{Tuần 5--6:}
\begin{itemize}
\item \textbf{Phú:} Dựng skeleton framework v0 (Knowledge + Workflow + Tools).
Prototype Admissions Agent trên dữ liệu giả của 1 trường.
\item \textbf{Hoàng:} Lên danh sách 20 trường ĐH/CĐ miền Trung + miền Nam có
khả năng là ICP. Viết pitch v1.
\item \textbf{Hảo:} Liên hệ warm intro với 5 trường có quan hệ sẵn.
\end{itemize}

\textbf{Tuần 7--8:}
\begin{itemize}
\item \textbf{Hoàng + Hảo:} Thực hiện 5 discovery call đầu tiên. Ghi chép.
\item \textbf{Phú:} Hoàn thiện prototype Admissions Agent, demo được cho khách.
\item \textbf{3 người:} Weekly review --- cập nhật ICP, cập nhật use case, cập nhật pricing.
\end{itemize}

\subsubsection*{Tháng 3 --- Ký design partner đầu tiên}

\textbf{Tuần 9--10:}
\begin{itemize}
\item Hoàn thành 10 discovery call. Phân tích --- 3 khách hàng có xác suất ký cao nhất.
\item Viết proposal cho 3 khách. Demo prototype với dữ liệu thật của họ.
\item Bắt đầu đàm phán hợp đồng pilot (50--150 triệu).
\end{itemize}

\textbf{Tuần 11--12:}
\begin{itemize}
\item Ký ít nhất \textbf{1 hợp đồng pilot trả phí}.
\item Kick-off dự án. Phú bắt đầu build phiên bản thật.
\item \textbf{Quarterly offsite --- nhìn lại 90 ngày, hiệu chỉnh playbook.}
\end{itemize}

\begin{keybox}[title=Mốc thành công của 90 ngày]
\textbf{1 hợp đồng pilot trả phí + 10 discovery call + framework v0 có demo được + 3 người vẫn muốn đi cùng nhau.}

Nếu cả 4 điều trên đạt được, công ty đã vượt qua điểm khó nhất ban đầu.
\end{keybox}

\subsection{Cam kết của 3 người}

\vspace{1em}

Với chữ ký dưới đây, chúng tôi --- \textbf{Đỗ Phúc Hảo, Lê Hữu Phú, Nguyễn Ngọc
Hoàng} --- đồng thuận:

\begin{enumerate}
\item Cam kết ít nhất \textbf{3 năm} đi cùng nhau trong công ty này.
\item Chấp nhận vai trò, quyền quyết định, và KPI đã nêu trong playbook.
\item Tuân thủ cơ chế \textbf{3 chữ ký} cho các quyết định lớn.
\item Không vi phạm 7 \textbf{red lines} đã liệt kê.
\item Cam kết \textbf{minh bạch hoàn toàn} về tài chính, lịch làm việc, và quyết định.
\item Đồng ý hiệu chỉnh playbook này sau mỗi \textbf{quarterly offsite}, khi cả 3 người đồng thuận.
\item Đối xử với nhau --- kể cả khi không đồng ý --- bằng \textbf{sự tôn trọng, kiên nhẫn, và giả định thiện chí}.
\end{enumerate}

\vspace{2em}

\begin{center}
\begin{tabular}{ccc}
\rule{5cm}{0.4pt} & \rule{5cm}{0.4pt} & \rule{5cm}{0.4pt} \\
\textbf{Đỗ Phúc Hảo} & \textbf{Lê Hữu Phú} & \textbf{Nguyễn Ngọc Hoàng} \\
CEO / Strategy & Platform Lead & GTM Lead \\
\small Ngày: \hspace{2cm} & \small Ngày: \hspace{2cm} & \small Ngày: \hspace{2cm} \\
\end{tabular}
\end{center}

\vfill

\begin{center}
\small\itshape\color{mid}
``Ba người đi cùng nhau, tất yếu có người làm thầy ta.'' --- \textit{Khổng Tử}\\[0.5em]
Nhưng trong công ty này, \textbf{chúng tôi đi cùng nhau và làm thầy lẫn nhau.}
\end{center}

\end{document}
